\chapter*{Introduction}
\addcontentsline{toc}{chapter}{Introduction}

\hspace{1cm} Ce second rapport du projet GE S8 présente la phase de conception détaillée d'une installation domestique autonome en énergie. Après le cadrage fonctionnel établi au rapport 1, l'objectif est ici de transformer les besoins en solutions techniques concrètes, justifiées et réalisables. Le document se concentre donc sur les études comparatives, les choix d'architecture, le dimensionnement des composants et la conception des cartes électroniques.

La démarche suivie repose sur un enchaînement clair : analyse des contraintes, comparaison de plusieurs topologies, sélection argumentée des composants, puis validation par calculs et simulations. Cette logique est appliquée à l'ensemble de la chaîne, depuis l'alimentation (conversion AC/DC et DC/DC) jusqu'à la distribution des usages et à la supervision (bus CAN, capteurs, commande et interface utilisateur). Les schémas électriques détaillés et le routage PCB associés traduisent ces choix en solutions matérielles prêtes à être fabriquées.

Le rapport est organisé en quatre volets techniques : une vue d'ensemble du système, la partie alimentation, la partie distribution et la partie supervision/interface. Chaque chapitre présente les hypothèses de travail, les critères de décision, les résultats de dimensionnement et les conceptions retenues, afin d'assurer la cohérence globale du système et de préparer la phase d'assemblage et de tests.